\section*{Introduction: Why This Chapter Exists Separately}

Chapter 15 developed willingness to pay as a property of the buyer, the job, and the alternatives. But a market may contain several kinds of buyers, multiple benefits, and many possible units of exchange. The founder must decide whether to sell per product, per seat, per year, per project, per outcome, or as a bundle. They must decide which work belongs in the basic offer, which service levels should cost more, how to price uncertainty, when to discount, and how to test a price without damaging trust. These are \textbf{pricing-design decisions}, not simply willingness-to-pay estimates.

This chapter is an extended development of the founder-supplied essay \emph{The Theory of Pricing and Value Capture}. Its central insights: including the value-capture triangle, 14 sources of defensible premiums, price waterfalls, purchase evidence, and Bayes Institute's proof-based premium: are incorporated and sharpened here. Some earlier essay material properly belongs in Chapter 15 (buyer psychology) and Chapter 17 (complete business economics), and so appears here only where it changes an actual pricing decision.

\paragraph{Method note.} Monetary examples, sample demand curves, conversion rates, cohort sizes, and package prices are hypothetical teaching devices; they are neither current market-price surveys nor Bayes Institute recommendations. The simplified formulae deliberately isolate effects. The founder must replace their assumptions with actual transactions, full costs, and local legal or accounting advice where appropriate.

\section*{Part I: Price as a Strategic Choice}

\subsection*{1. The Price Tag Is the Last Visible Component of a Much Larger Design}

2 businesses can deliver an essentially similar technical result and earn radically different economics. 1 sells hours of training to individuals; another sells an assessed capability program to an enterprise; a 3rd licenses a curriculum to colleges; a 4th publishes material free and sells intensive expert feedback. They may reuse some intellectual content, but their prices are not simply different figures on the same brochure. They represent different transactions with different buyers, service obligations, acquisition paths, and risks.

The common founder error is to finish the product, look at competitors, pick a fee near the middle, and then ask marketing to justify it. This forces a late numerical choice onto an earlier architectural decision. If the intended buyer values reliability rather than content volume, the seller needs to design reliability into the product and into the commercial promise. If the buyer values predictable spend, usage-based fees may fail even when they maximize the seller's modeled revenue. If expert attention is scarce, a flat ``unlimited support'' promise can make high demand financially dangerous.

\textbf{Pricing is the design of what is exchanged, in which units, under which conditions, between which parties, with which risks allocated to each.} The number comes last, though pricing feasibility should influence product design from the beginning.

\subsection*{2. 4 Tests Before Setting a Number}

A defensible price must pass 4 tests:
\begin{enumerate}
  \item \textbf{Customer value:} the buyer has a credible reason to choose the offer over feasible alternatives.
  \item \textbf{Economic viability:} the seller can acquire, deliver, support, and improve it at acceptable cost.
  \item \textbf{Transaction feasibility:} the price fits budgets, approval processes, timing, and payment risk.
  \item \textbf{Strategic legitimacy:} the terms align with desired reputation, access policy, regulation, and repeated trust.
\end{enumerate}

These tests constrain 1 another. A ₹1,000 product can sell easily but fail because acquisition costs ₹1,200. A ₹1 lakh program may have excellent potential margin but no credible way to demonstrate its educational value before purchase. A high-volume promotion might fill every seat and destroy feedback quality. An enterprise deal can have large accounting profit but such late cash collection that the founder cannot finance delivery.

Do not confuse a price test with proof of the entire business. The price tests demand for a \emph{specified proposition} under specified conditions. Changing price, packaging, promises, or target customer changes the object being tested.

\subsection*{3. The Value-Capture Triangle: Buyer, Alternative, Seller}

Think of every commercial offer as a triangle. At 1 corner sits the \textbf{buyer}, who values particular consequences. At the 2nd sits the \textbf{alternative}, which establishes the comparison, including doing nothing. At the 3rd sits the \textbf{seller}, whose cost structure determines the economics of supplying the difference. A price must make sense from all 3 corners.

Suppose a buyer values a service at ₹50,000 relative to no service. A substitute provides almost the same outcome for ₹20,000. The new seller's relevant premium is not automatically ₹30,000. They must establish what \emph{incremental} benefit their version provides and how much value the buyer can keep. On the other side, if delivering the differentiated version costs ₹28,000, an apparent customer opportunity at ₹25,000 is not commercially attractive under ordinary contribution logic.

The triangle guards against the opposite errors of cost-plus absolutism and pure willingness-to-pay extraction. The right business creates surplus on both sides of the exchange and can demonstrate why it deserves its share.
